% This is part of Mes notes de mathématique
% Copyright (c) 2011-2020
%   Laurent Claessens
% See the file fdl-1.3.txt for copying conditions.

%+++++++++++++++++++++++++++++++++++++++++++++++++++++++++++++++++++++++++++++++++++++++++++++++++++++++++++++++++++++++++++
\section{Module sur un anneau}
%+++++++++++++++++++++++++++++++++++++++++++++++++++++++++++++++++++++++++++++++++++++++++++++++++++++++++++++++++++++++++++

\begin{definition}[module sur un anneau\cite{ooJGVOooSjQBVh}]       \label{DEFooHXITooBFvzrR}
    Soit un anneau \( A\). Un \defe{module à gauche}{module!à gauche} sur \( A\) est la donnée d'un triple \( (M,+,\cdot)\) où
    \begin{enumerate}
        \item
            \( +\) est une loi de composition interne à \( M\), c'est-à-dire \( +\colon M\times M\to M\),
        \item
            \( \cdot\) est une loi de composition externe, c'est-à-dire \( \cdot\colon A\times M\to M\)
    \end{enumerate}
    telles que
    \begin{enumerate}
        \item
            \( (M,+)\) est un groupe\footnote{Nous verrons dans la proposition~\ref{PROPooGARGooDiMqtN} qu'il est forcément commutatif.}.
        \item
            \( a\cdot(x+y)=a\cdot x+a\cdot y\),
        \item
            \( (a+b)\cdot x=a\cdot x+b\cdot x\),
        \item
            \( (ab)\cdot x=a\cdot(b\cdot x)\)
        \item
            \( 1\cdot x=x\).
    \end{enumerate}
    pour tout \( a,b\in A\) et \( x,y\in M\).
\end{definition}

\begin{proposition}\label{PROPooGARGooDiMqtN}
    Si \( M\) est un module sur un anneau, alors \( (M,+)\) est un groupe commutatif.
\end{proposition}

\begin{proof}
    Il suffit de calculer \( (1+1)\cdot (x+y)\) de deux façons différentes :
    \begin{equation}
        (1+1)\cdot (x+y)=1\cdot (x+y)+1\cdot (x+y)=x+y+x+y
    \end{equation}
    d'une part et
    \begin{equation}
        (1+1)\cdot (x+y)=(1+1)\cdot x+(1+1)\cdot y=x+x+y+y,
    \end{equation}
    d'autre part. En égalant les deux expressions, il vient
    \begin{equation}
        x+y+x+y=x+x+y+y,
    \end{equation}
    qui se simplifie (nous sommes dans un groupe) en \( y+x=x+y\).
\end{proof}

\begin{definition}\label{DEFooKHWZooIfxdNc}
    Un \defe{espace vectoriel}{espace!vectoriel} est un module sur un corps commutatif\footnote{La condition de commutativité n'est pas indispensable, mais comme nous ne parlerons que de corps commutatifs\ldots}.
\end{definition}

\begin{definition}[\cite{BIBooSTWWooItiMUp}]        \label{DEFooRUKVooLnXxdS}
    Soient un \( A\)-module \( M\) et un ensemble \( I\). Une famille \( (m_i)_{i\in I}\) est \defe{libre}{partie libre!module} si ils sont \defe{linéairement indépendants}{linéairement indépendant!module}, c'est-à-dire si pour tout choix d'une partie finie \( J\) dans \( I\) et d'éléments \( (a_j)_{j\in J}\) dans \( A\), si nous avons
    \begin{equation}
        \sum_{j\in J}a_jm_j=0,
    \end{equation}
    alors \( a_j=0\) pour tout \( j\).
\end{definition}

\begin{definition}[\cite{BIBooNKWVooYfrwSd}]        \label{DEFooWBOBooJNyyBF}
    Soit \( S\), une partie du \( A\)-module \( M\). Le \defe{sous-module engendré}{sous-module engendré} par \( S\) est l'ensemble des éléments de \( M\) qui sont des combinaisons linéaires finies d'éléments de \( S\), c'est-à-dire de sommes de la forme
    \begin{equation}
        \sum_{t\in T}a_tt
    \end{equation}
    où \( T\) est fini dans \( S\) et \( a_t\in A\).
\end{definition}

%--------------------------------------------------------------------------------------------------------------------------- 
\subsection{Module produit}
%---------------------------------------------------------------------------------------------------------------------------

\begin{lemmaDef}[\cite{BIBooSTWWooItiMUp}]        \label{DEFooLCJEooBvVmkV}
    Soient un anneau \( A\) et un ensemble \( I\). Le \( A\)-\defe{module produit}{module produit} \( A^I\) est l'ensemble des applications \( I\to A\).

    En termes de notations, nous écrivons ceci :
    \begin{equation}
        A^I=\{ (a_i)_{i\in I},a_i\in A \}.
    \end{equation}
    L'ensemble \( A^I\) devient un module par les définition, pour \( x,y\in A^I\) et \( a\in A\) :
    \begin{subequations}
        \begin{align}
            ax&=(ax_i)_{i\in I}\\
            x+y&=(x_i+y_i)_{i\in I}     \label{EQooODBMooQKLUgd}.
        \end{align}
    \end{subequations}
    En d'autres termes, \( A^I=\Fun(I,A)\).
\end{lemmaDef}

\begin{lemma}
    Pour chaque \( i\in I\) nous considérons l'élément \( e_i\in A^I\) donné par
    \begin{equation}
        e_i=(\delta_{ij})_{j\in I}.
    \end{equation}
    La famille \( \{ e_i \}_{i\in I}\) est libre\footnote{Définition \ref{DEFooRUKVooLnXxdS}.} dans \( A^I\).
\end{lemma}

\begin{proof}
    Soient \( J\) fini dans \( I\) ainsi que des éléments \( a_j\in A\) (\( j\in J\)). Nous supposons que\footnote{Pour rappel, les sommes finies sont définies par \ref{DEFooLNEXooYMQjRo}.} \( \sum_{j\in J}a_je_j=0\). Calculons un peu :
    \begin{equation}
        \sum_{j\in J}a_je_j=\sum_{j\in J}(a_j\delta_{ji})_{i\in I}=\left( \sum_{j\in J}a_j\delta_{ji} \right)_{i\in I}.
    \end{equation}
    Pour que le tout soit nul dans \( A^I\), il faut que
    \begin{equation}
        \sum_{j\in J}a_j\delta_{ji}
    \end{equation}
    soit nul pour tout \( i\in I\). Si nous fixons \( i\in I\), la somme sur \( j\) possède un seul terme non annulé par \( \delta_{ji}\), et c'est le terme \( j=i\). Nous avons donc \( a_i=0\).
\end{proof}

\begin{definition}      \label{DEFooBMEPooFsCHgb}
    Nous notons \( A^{(I)}\) le sous-module de \( A^I\) engendré\footnote{Définition \ref{DEFooWBOBooJNyyBF}.} par les \( e_i\).
\end{definition}

\begin{theorem}[Propriété universelle de \( A^{(I)}\)\cite{BIBooSTWWooItiMUp}]      \label{THOooPDZCooJnHbOd}
    Soient un anneau \( A\) ainsi qu'un \( A\)-module \( P\). Pour \( \phi\in\Hom_A(A^{(I)}, P)\), nous considérons
    \begin{equation}
        \begin{aligned}
            \phi|_I\colon I&\to P \\
            i&\mapsto \phi(e_i). 
        \end{aligned}
    \end{equation}
    \begin{enumerate}
        \item
            
    L'application
    \begin{equation}
        \begin{aligned}
            f\colon \Hom_A(A^{(I)},P)&\to \Fun(I,P) \\
            \phi&\mapsto \phi|_I 
        \end{aligned}
    \end{equation}
    est une bijection.
\item
    L'application inverse est \( g\colon \Fun(I,P)\to \Hom_A(A^{(I)},P) \) donnée par
    \begin{equation}
        g(\psi)\big( \sum_{j\in J}a_je_j \big)=\sum_{j\in J}a_j\psi(j)
    \end{equation}
    pour tout \( J\) fini dans \( I\) et choix de \( a_j\in A\).
    \end{enumerate}
\end{theorem}

\begin{proof}
    Nous allons montrer que \( g\big( f(\phi) \big)=\phi\) et que \( f\big( g(\psi) \big)=\psi\) pour tout \( \phi\in\Hom_A(A^{(I)},P)\) et pour tout \( \psi\in \Fun(I,P)\).

    Dans un premier sens nous avons :
    \begin{subequations}
        \begin{align}
            g\big( f(\phi) \big)\big( \sum_ja_je_j \big)&=\sum_ja_jf(\phi)(j)\\
            &=\sum_ja_j\phi(e_j)\label{SUBALIGNooBWPLooHeIaQg}\\
            &=\phi(\sum_ja_je_j)        \label{SUBALIGNooUOQPooCwLgZo}.
        \end{align}
    \end{subequations}
    Justifications :
    \begin{itemize}
        \item 
            Pour \eqref{SUBALIGNooBWPLooHeIaQg}, nous avons utilisé le fait que \( f(\phi)(i)=\phi|_I(i)=\phi(e_i)\).
        \item
            Pour \eqref{SUBALIGNooUOQPooCwLgZo}, nous utilisons le fait que \( \phi\) est un morphisme de modules.
    \end{itemize}
    Et pour l'autre sens,
    \begin{equation}
        f\big( g(\psi) \big)(i)=g(\psi)(e_i)=\psi(i).
    \end{equation}
\end{proof}

%--------------------------------------------------------------------------------------------------------------------------- 
\subsection{Sous-module}
%---------------------------------------------------------------------------------------------------------------------------

Soient \( M\) un \( A\)-module et \( x=(x_i)_{i\in I}\) une famille d'éléments de \( M\) paramétrée par l'ensemble \( I\). Nous considérons l'application
\begin{equation}
    \begin{aligned}
        \mu_x\colon A^{(I)}&\to M \\
        (a_i)_{i\in I}&\mapsto \sum_{i\in I}a_ix_i.
    \end{aligned}
\end{equation}
Ici \( A^{(I)}\) désigne l'ensemble de toutes les applications \( I\to A\) de support fini.

\begin{definition}      \label{DefBasePouyKj}
    À l'instar des espaces vectoriels, les modules ont une notion de partie libre, génératrice et de bases :
    \begin{enumerate}
        \item
            Si \( \mu_x\) est surjective, nous disons que \( x\) est une partie \defe{génératrice}{génératrice!partie d'un module}.
        \item
            Si \( \mu_x\) est injective, nous disons que la partie \( x\) est \defe{libre}{libre!partie d'un module}.
        \item
            Si \( \mu_x\) est bijective, nous disons que la partie \( x\) est une \defe{base}{base!d'un module}.
    \end{enumerate}
\end{definition}

\begin{definition}
  Un sous-ensemble \( N\subset M\) est un \defe{sous-module}{sous-module} si \( (N,+)\) est un sous-groupe de \( (M,+)\) et si \( a\cdot x\in N\) pour tout \( x\in N\) et pour tout \( a\in A\).
\end{definition}

\begin{example}
    Un anneau \( A\) est lui-même un \( A\)-module et ses sous-modules sont les idéaux.
\end{example}

\begin{definition}
    Soit \( M\) un module sur un anneau commutatif \( A\). Un \defe{projecteur}{projecteur!dans un module} est une application linéaire \( p\colon M\to M\) telle que \( p^2=p\).

    Une famille \( (p_i)_{i\in I}\) sur \( M\) est \defe{orthogonale}{orthogonal!famille de projecteurs} si \( p_i\circ p_j=0\) pour tout \( i\neq j\). La famille est \defe{complète}{complète!famille de projecteurs} si \( \sum_{i\in I}p_i=\mtu\).
\end{definition}

\begin{theorem}     \label{ThoProjModpAlsUR}
    Soient des sous modules \( M_1,\ldots,M_n\) du module \( M \) tels que \( M=M_1\oplus\ldots\oplus M_n\). Les applications \( p_i\) définies par
    \begin{equation}
        p_i(x_1+\ldots+x_n)=x_i
    \end{equation}
    forment une famille orthogonale de projecteurs et \( p_1+\cdots +p_n=\id\).

    Inversement, si \( (p_1,\ldots, p_n)\) est une famille orthogonale de projecteurs dans un module \( \modE\) tel que \( \sum_{i=1}^np_i=\id\), alors
    \begin{equation}
        M=\bigoplus_{i=1}^np_i(M).
    \end{equation}
\end{theorem}

\begin{definition}
    Un module est \defe{simple}{simple!module}\index{module!simple} ou \defe{irréductible}{irréductible!module}\index{module!irréductible} s'il n'a pas d'autres sous-modules que \( \{ 0 \}\) et lui-même. Un module est \defe{indécomposable}{indécomposable!module}\index{module!indécomposable} s'il ne peut pas être écrit comme somme directe de sous-modules.
\end{definition}

Un module simple est a fortiori indécomposable. L'inverse n'est pas vrai comme le montre l'exemple suivant.

\begin{example}
    Soit \( \modE=\eC[X]/(X^2)\) vu comme \( \eC[X]\)-module. C'est le \( \eC[X]\)-module des polynômes de la forme \( aX+b\) avec \( a,b\in \eC\). L'ensemble des polynômes de la forme \( aX\) est un sous-module. Le module \( \modE\) n'est donc pas simple. Il est cependant indécomposable parce que \( \{ aX \}\) est le seul sous-module non trivial. En effet si \( \modF\) est un sous-module de \( \modE\) contenant \( aX+b\) avec \( b\neq 0\), alors \( \modF\) contient \( X(aX+b)=bX\) et donc contient tout \( \modE\).
\end{example}

\begin{definition}[Algèbre\cite{ZSyHmiy}]   \label{DefAEbnJqI}
    Si \( \eK\) est un corps commutatif\footnote{Définition~\ref{DefTMNooKXHUd}}, une \( \eK\)-\defe{algèbre}{algèbre} \( A\) est un espace vectoriel\footnote{Définition~\ref{DEFooKHWZooIfxdNc}.} muni d'une opération bilinéaire \( \times\colon A\times A\to A\), c'est-à-dire telle que pour tout \( x,y,z\in A\) et pour tout \( \alpha,\beta\in\eK\),
    \begin{enumerate}
        \item
            \( (x+y)\times z=x\times z+y\times z\)
        \item
            \( x\times (y+z)=x\times y+x\times z\)
        \item
            \( (\alpha x)\times (\beta y)=(\alpha\beta)(x\times y)\).
    \end{enumerate}
    Si \( A\) et \( B\) sont deux \( \eK\)-algèbres, une application \( f\colon A\to B\) est un \defe{morphisme d'algèbres}{morphisme!d'algèbres} entre \( A\) et \( B\) si pour tout \( x,y\in A\) et pour tout \( \alpha\in \eK\),
    \begin{enumerate}
        \item
            \( f(xy)=f(x)f(y)\)
        \item
            \( f(x+\alpha y)=f(x)+\alpha f(y)\)
    \end{enumerate}
    où nous avons noté \( xy\) pour \( x\times y\).
\end{definition}

\begin{lemma}[\cite{MonCerveau}]   \label{LEMooVKLKooSAHmpZ}
    Soient une algèbre \( A\) et une famille \( (X_i)_{i\in I}\) de sous-algèbres de \( A\) (ici \( I\) est un ensemble quelconque). Alors la partie \( X=\bigcap_{i\in I}X_i\) est une sous-algèbre de \( A\).
\end{lemma}

\begin{proof}
    Nous devons prouver que si \( x\) et \( y\) sont dans \( X\) et \( \lambda\in \eK\), alors \( xy\), \( x+y\) et \( \lambda x\) sont dans \( X\). Pour tout \( i\in I\) nous avons \( x,y\in X_i\) et donc \( xy\in X_i\), \( x+y\in X_i\) et \( \lambda x\in X_i\) (parce que \( X_i\) est une algèbre). Donc \( xy\),\( x+y\) et \( \lambda x\) sont dans \( X_i\) pour tout \( I\), et donc dans \( X\).
\end{proof}

\begin{definition}\label{DefkAXaWY}
    L'\defe{algèbre engendrée}{algèbre!engendrée} par \( X\) est l'intersection de toutes les sous-algèbres de \( A\) contenant \( X\) (qui est une algèbre par le lemme~\ref{LEMooVKLKooSAHmpZ}).
\end{definition}

%+++++++++++++++++++++++++++++++++++++++++++++++++++++++++++++++++++++++++++++++++++++++++++++++++++++++++++++++++++++++++++
\section{Polynômes}
%+++++++++++++++++++++++++++++++++++++++++++++++++++++++++++++++++++++++++++++++++++++++++++++++++++++++++++++++++++++++++++
\label{SEC_annPolynomes}

%--------------------------------------------------------------------------------------------------------------------------- 
\subsection{Polynômes d'une variable}
%---------------------------------------------------------------------------------------------------------------------------

Et voilà la définition que tout le monde attendait; la définition des anneaux de polynômes. Pour ne pas taper trop fort du premier coup, nous commençons par les polynômes d'une seule variable.

Nous allons définir et étudier ici l'anneau des polynômes sur un anneau \( A\), c'est-à-dire ce qui sera noté \( A[X]\). Pour \( \eK(X)\) lorsque \( \eK\) est un corps, voir~\ref{DEFooQPZIooQYiNVh}.

L'ensemble des polynômes sur \( A\) sera simplement \( A^{(\eN)}\). Vu que \( \eN\) est un ensemble bien particulier possédant plein de structure, nous allons pouvoir mettre sur \( A^{(\eN)}\) une structure non seulement de \( A\)-module (ça c'est déjà fait), mais en plus d'anneau ainsi qu'une évaluation.
\begin{definition}      \label{DEFooFYZRooMikwEL}
    L'ensemble des \defe{polynômes}{polynômes} en une indéterminée sur l'anneau \( A\) est l'anneau \( A^{(\eN)}\) que nous avons défini en \ref{DEFooBMEPooFsCHgb}.
\end{definition}

Nous ne voyons pas encore bien le lien avec la notation \( A[X]\); nous le verrons plus tard. 

Vu que \( A^{(\eN)}\) est engendré par les \( e_i\), tout polynôme sur \( A\) s'écrit \( P=\sum_{i=1}^na_ie_i\).

\begin{definition}      \label{DEFooNXKUooLrGeuh}
    Nous ajoutons deux structures à \( A^{(\eN)}\).
    \begin{description}
        \item[L'évaluation] Si \( \alpha\in A\) et si \( P\in A^{(\eN)}\), nous définissons \( P(\alpha)\) par
            \begin{equation}        \label{EQooDJISooTEkMOw}
                P(\alpha)=(\sum_{i=0}^{n}a_ie_i)(\alpha)=\sum_{i=0}^na_i\alpha^i,
            \end{equation}
            étant entendu que \( \alpha^0=1\) dans \( A\).

            Cette définition s'étend immédiatement au cas où \( B\) est un anneau qui étend \( A\). Dans ce cas nous pouvons définir \( P(b)\) pour tout \( P\in \eA^{(\eN)}\) et \( b\in B\) avec la même formule \eqref{EQooDJISooTEkMOw}.
        \item[Le produit] C'est ici que la structure particulière de \( \eN\) est utilisée. Nous définissons le produit \( A^{\eN}\times A^{(\eN)}\to A^{(\eN)}\) de la façon suivante. Si \( (P_k)_{k\in \eN}\) est la suite (presque partout nulle) d'éléments de \( A\) qui définit \( P\) et si \( (Q_k)_{k\in \eN}\) est celle de \( Q\), nous notons
        \begin{equation}    \label{EQooTNCSooKklisb}
            (PQ)_n=\sum_{k=0}^nP_kQ_{n-k},
        \end{equation}
        et donc \( PQ=\sum_i(PQ)_ie_i\). Plus explicitement,
        \begin{equation}    \label{EQooCIBUooAgpxjL}
            (\sum_{i=0}^na_ie_i)(\sum_{j=0}^mb_je_j)=\sum_{k=0}^{\infty}\Big( \sum_{\substack{  (i,j)\in \eN^2 \\i+j=k}}a_ib_j \Big)e_k.
        \end{equation}
        Notons qu'à droite, la somme sur \( k\) est une somme finie.
    \end{description}
\end{definition}

\begin{proposition}     \label{PROPooGDQCooHziCPH}
    Soit un anneau \( A\). À propos de structure sur \( A^{(\eN)}\).
    \begin{enumerate}
        \item
            Avec le produit, l'ensemble \( A^{(\eN)}\) devient un anneau.
        \item
    L'application
    \begin{equation}
        \begin{aligned}
            g\colon A^{(\eN)}&\to A \\
            P&\mapsto P(\alpha)
        \end{aligned}
    \end{equation}
    est un morphisme d'anneaux\footnote{Définition \ref{DEFooSPHPooCwjzuz}.}. En particulier, \( (PQ)(\alpha)=P(\alpha)Q(\alpha)\).
    \end{enumerate}
\end{proposition}

\begin{proof}
    En plusieurs points
    \begin{subproof}
        \item[Anneau]
            L'identité pour le produit dans \( A^{(\eN)}\) est le polynôme donné par \( a_0=1\) et \( a_i=0\) pour \( i\neq 0\). Cela se vérifie en utilisant directement la définition \eqref{EQooCIBUooAgpxjL}. La distributivité aussi\quext{Je n'ai pas fait les calculs, écrivez-moi pour me dire si ça va facilement.}.
        \item[Le morphisme]
    Nous notons \( P_k\) les éléments de la suite définissant \( P\) et \( Q_k\) ceux de \( Q\). Alors nous avons
    \begin{equation}
        (P+Q)(\alpha)=\sum_k(P_k+Q_k)\alpha^k=\sum_kP_k\alpha^k+\sum_kQ_k\alpha^k=P(\alpha)+Q(\alpha).
    \end{equation}
    Vous aurez noté que la première égalité était la définition \eqref{EQooODBMooQKLUgd}. De même,
    \begin{subequations}
        \begin{align}
            P(\alpha)Q(\alpha)&=\big( \sum_nP_n\alpha^n \big)\big( \sum_kQ_k\alpha^k \big)=\sum_kQ_k\big( \sum_nP_n\alpha^n \big)\alpha^k=\sum_k\sum_nQ_kP_n\alpha^{n+k}\\
            &=\sum_m\big( \sum_{l=0}^mP_lQ_{m-l} \big)\alpha^m=\sum_m(PQ)_m\alpha^m=(PQ)(\alpha).
        \end{align}
    \end{subequations}
    \end{subproof}
\end{proof}

%--------------------------------------------------------------------------------------------------------------------------- 
\subsection{La notation \texorpdfstring{$ A[X]$}{A[X]}}
%---------------------------------------------------------------------------------------------------------------------------
\label{SUBSECooLEKVooFBPSJz}

Si \( A\) est un anneau, nous avons déjà défini les polynômes en une indéterminée sur \( A\) comme étant le module \( A^{(\eN)}\) qui est devenu un anneau par la proposition \ref{PROPooGDQCooHziCPH}.

Le polynôme donné par la suite \( (a_n)_{n\in \eN}\) est souvent notée
\begin{equation}
    \sum_ka_kX^k.
\end{equation}
Par exemple avec \( a=(4,2,8,0, \dots)\) nous avons \( a=8X^2+2X+4\). Nous utiliserons souvent cette notation, qui est très pratique parce qu'elle s'adapte bien aux règles de multiplication et d'addition, en particulier la distributivité.

Il y a (au moins) deux façons de comprendre ce que signifie réellement «\( X\)» dans cette notation.

%///////////////////////////////////////////////////////////////////////////////////////////////////////////////////////////
\subsubsection{Première façon (qui botte en touche)}
%///////////////////////////////////////////////////////////////////////////////////////////////////////////////////////////

La première est de dire qu'il n'a pas de significations, et que \( X^2\) est un simple abus de notations pour écrire \( (0,0,1,0,\cdots)\). Avec cette façon de voir, nous notons l'anneau des polynômes sur \( A\) par «\( A[X]\)» où le \( X\) n'a pas d'autres raisons d'être que d'avertir le lecteur que nous réservons la lettre «\( X\)» pour utiliser la notation pratique des polynômes.

%///////////////////////////////////////////////////////////////////////////////////////////////////////////////////////////
\subsubsection{Seconde façon (la bonne)}
%///////////////////////////////////////////////////////////////////////////////////////////////////////////////////////////
\label{SUBSUBSECooPNBYooWXEHrg}

La seconde façon de voir le «\( X\)» est de nous rappeler que \( A^{(\eN)}\) a une base en tant que module (les \( e_k\) dont nous avons parlé plus haut), et que nous l'avons muni d'une structure multiplicative. Nous posons \( X=e_1\), et nous prenons la convention \( X^0=1\). Il s'ensuit alors ceci:

\begin{propositionDef}
  L'anneau \(A^{(\eN)}\), muni de l'addition "composante par composante" (comme \( A \)-module) et du produit défini à l'aide de~\eqref{EQooTNCSooKklisb}, est l'anneau de polynômes \( A[X]\).
\end{propositionDef}

Cela n'est pas si évident, car il faut bien s'assurer que "\( e_k=X^k\)".

Dans les deux cas, il n'est a priori pas légitime d'écrire des égalités comme « \( P(X)=X^2+2X-3\) », et encore moins de dire «Le polynôme \( P\), \emph{évalué} en \( X\) vaut \( X^2+2X-3\)»  : il est plus correct d'écrire « \( P=X^2+2X-3\) ».

Néanmoins, le lemme suivant montre que ces notations tombent vraiment à point. La véritable difficulté de l'énoncé est de comprendre qu'il n'est pas trivial.

Nous avons vu dans la définition~\ref{DEFooNXKUooLrGeuh} que si \( B\) est un anneau qui étend \( A\), et si \(P\in A[X] \), alors nous avons une définition de \( P(b)\) pour tout \( b\in B\). Nous appliquons cela à \( B=A[X]\), qui est un anneau qui étend \( A\). Autrement dit, si \( P\) et \( Q\) sont des polynômes, ça a un sens d'écrire \( P(Q)\) et le résultat sera un élément de \( A[X]\). 

Dans le cas particulier \( Q=X\), nous avons une chouette formule.
\begin{lemma}       \label{LEMooGKWQooVOyeDX}
    Nous avons
    \begin{equation}
        P(X)=P
    \end{equation}
    pour tout \( P\in A[X]\).
\end{lemma}

\begin{proof}
    Si \( P=(a_k)_{k\in \eN}\) alors par définition \( P(\alpha)=\sum_ka_k\alpha^k\) dès que \( \alpha\) est dans un anneau \( B\) qui étend \( A\). Nous considérons le cas particulier \( B=\eA[X]\) et \( \alpha=X\), c'est-à-dire \( Q=(0,1,0,\ldots)\), l'élément \( P(X)\) de \( A[X]\) vaut
    \begin{equation}        \label{EQooABULooFCEasf}
        \sum_ka_kX^k,
    \end{equation}
    qui est exactement \( P\) lui-même.
\end{proof}

Mais il faut bien comprendre que si \( P\) est le polynôme \( (-3,2,1,0,\ldots)\), noté \( X^2+2X-3\), écrire \( P(X)=X^2+2X-3\) est une pirouette de notations que rien ne justifie par rapport à simplement écrire \( P=X^2+2X-3\).

%---------------------------------------------------------------------------------------------------------------------------
\subsection{Polynômes de plusieurs variables}
%---------------------------------------------------------------------------------------------------------------------------

\begin{definition}      \label{DEFooZNHOooCruuwI}
    L'ensemble des \defe{polynôme de \( n\) variables}{polynôme de plusieurs variables} sur l'anneau \( A\) est \( A^{(\eN^n)}\), c'est-à-dire l'ensemble des suites indexées par \( \eN^n\) et dont seulement une quantité finie de coefficients sont non nuls.

    Le produit sur \( A[X_1,\ldots, X_n]\) est défini par
    \begin{equation}
        (PQ)(k_1,\ldots, k_n)=\sum_{\substack{ (l_1,\ldots, l_n),(m_1,\ldots, m_n)\in \eN^n\times \eN^n   \\l_i+m_i=k_i}}P_{l_1,\ldots, l_n}Q_{m_1,\ldots, m_n}.
        \end{equation}
\end{definition}

\begin{normaltext}
    Dans \( A[X_1,\ldots, X_n]\), la multiplication n'est pas la multiplication de fonctions \( \eN^n\to \eK\), parce que le but est d'obtenir la multiplication usuelle au niveau des évaluations.
\end{normaltext}

\begin{definition}
    Si \( P\) est un polynôme de \( n\) variables sur \( A\), et si \( (x_1,\ldots, x_n)\in A^n\), \defe{l'évaluation}{évaluation!polynôme plusieurs variables} de \( P\) sur \( (x_1,\ldots, x_n)\) est
    \begin{equation}
        P(x_1,\ldots, x_n)=\sum_{(k_1,\ldots, k_n)\in \eN^n}P_{k_1,\ldots, k_n}x_1^{k_1}\ldots x_n^{k_n}.
    \end{equation}
    Notez que la somme, bien que sur \( \eN^n\), est une somme finie.
\end{definition}

\begin{normaltext}
    Comme dans le cas des polynômes d'une seule variable, les \( X_i\) dans la notation \( A[X_1,\ldots, X_n]\) sont à prendre à la légère. L'anneau des polynômes de \( n\) variables sur \( A\) aurait mieux fait d'être noté par exemple par \( \mP_n(A)\).

    Le fait est que nous avons les polynômes élémentaires définis par
    \begin{equation}
        X_1(k_1,\ldots, k_n)=\begin{cases}
            1    &   \text{si } (k_1,\ldots, k_n)=(1,0\ldots, 0)\\
            0    &    \text{sinon. }
        \end{cases}
    \end{equation}
    et que l'anneau des polynômes peut être vu comme \( A\) (les polynômes constants) étendus par les \( X_i\).

    Quoi qu'il en soit, les \( X_1\) dans la notation \( A[X_1,\ldots, X_n]\) sont des indices muets. L'anneau \( A[X_1,\ldots, X_n]\) est exactement le même que \( A[T_1,\ldots, T_n]\).
\end{normaltext}

%--------------------------------------------------------------------------------------------------------------------------- 
\subsection{Action du groupe symétrique}
%---------------------------------------------------------------------------------------------------------------------------

Par souci de notations, nous notons \( \Poly_n(A)\) l'anneau des polynômes de \( n\) variables sur \( A\). La propriété universelle de \( \Poly_n(A)=A^{(\eN^n)}\) du théorème \ref{THOooPDZCooJnHbOd} nous donne une application
\begin{equation}
    g\colon \Fun\big(\eN^n,\Poly_n(A)\big)\to \Hom_A\big( \Poly_n(A),\Poly_n(A) \big)
\end{equation}
Avec cela nous pouvons énoncer et démontrer le lemme qui donne l'action de \( S_n\)\footnote{Définition du groupe symétrique \( S_n\) en \ref{DEFooJNPIooMuzIXd}.} sur \( \Poly_n(A)\).

\begin{lemma}[\cite{BIBooFDZDooJQLjlB}]       \label{LEMooIRVQooHvoNBq}
    Pour \( \sigma\in S_n\) nous définissons 
    \begin{equation}
        \begin{aligned}
            \phi_{\sigma}\colon \eN^n&\to \Poly_n(A) \\
            m&\mapsto e_{\sigma(m)}. 
        \end{aligned}
    \end{equation}
    Alors l'application
    \begin{equation}
        \begin{aligned}
            \rho\colon S_n&\to \Hom_A\big( \Poly_n(A),\Poly_n(A) \big) \\
            \sigma&\mapsto g(\phi{\sigma}) 
        \end{aligned}
    \end{equation}
    est une action\footnote{Définition \ref{DefActionGroupe}.}.
\end{lemma}

\begin{proof}
    Nous commençons par donner une expression à notre \( \rho\). Un élément de \( \Poly_n(A)\) est de la forme \( \sum_{m\in \eN^n}a_me_m\), et nous avons\footnote{La somme est définie par \ref{DEFooLNEXooYMQjRo}, et ça va être important. Ah oui, en réalité partout, les sommes sont finies parce que les \( a_m\) (\( m\in \eN^n\)) sont presque tous nuls. Il faudrait écrire sur la somme sur \(\{ m\in \eN^2\tq a_m\neq 0 \}\), mais vous vous imaginez la complication dans la notation.}
    \begin{equation}
        \rho(\sigma)\big( \sum_{m\in \eN^n}a_me_m \big)=\sum_ma_m\phi_{\sigma}(m)=\sum_ma_me_{\sigma(m)}.
    \end{equation}
    
    Nous avons tout de suite \( \rho(\id)=\id\).

    En ce qui concerne la composition, nous avons d'une part
    \begin{equation}
        \rho(\sigma_1\sigma_2)\big( \sum_ma_me_m \big)=g(\phi_{\sigma_1\sigma_2})\big( \sum_ma_me_m \big)=\sum_ma_me_{\sigma_1\sigma_2(m)},
    \end{equation}
    et d'autre part,
    \begin{subequations}
        \begin{align}
            \rho(\sigma_1)\rho(\sigma_2)\big( \sum_ma_me_m \big)&=\rho(\sigma_1)\big( \sum_ma_me_{\sigma_2(m)} \big)\\
            &=\rho(\sigma_1)\big( \sum_ma_{\sigma_2^{-1}(m)}e_m \big)   \label{SUBEQooTSCYooCUWiRz}\\
            &=\sum_ma_{\sigma_2^{-1}(m)}e_{\sigma_1(m)}\\
            &=\sum_ma_me_{\sigma_1\sigma_2(m)}      \label{SUBEQooQPGPooVvqJdT}
        \end{align}
    \end{subequations}
    La proposition \ref{PROPooJBQVooNqWErk} est utilisée pour \eqref{SUBEQooTSCYooCUWiRz} et pour \eqref{SUBEQooQPGPooVvqJdT}.
\end{proof}

%+++++++++++++++++++++++++++++++++++++++++++++++++++++++++++++++++++++++++++++++++++++++++++++++++++++++++++++++++++++++++++
\section{Polynômes à coefficients dans un anneau commutatif}
%+++++++++++++++++++++++++++++++++++++++++++++++++++++++++++++++++++++++++++++++++++++++++++++++++++++++++++++++++++++++++++
\label{SECooVMABooVdhbPo}

\begin{lemma}       \label{LEMooXFMAooMBgIrN}
    Nous considérons un polynôme \( P\in A[X]\), et le quotient \( A[X]/(P)\). Pour tout polynôme \( Q\in A[X]\) nous avons les égalités
    \begin{equation}
        Q(\bar X)=\overline{ Q(X) }=\bar Q.
    \end{equation}
\end{lemma}

\begin{proof}
    Si \( Q=\sum_ka_kX^k\), alors par la linéarité de la prise de classes,
    \begin{equation}        \label{EQooXQRMooIPGFVM}
        \bar Q=\sum_ka_k\overline{ X^k }.
    \end{equation}
    Nous insistons sur le fait que cette égalité n'est rien d'autre que l'itération de la définition de la somme dans l'espace quotient : \( \bar x+\bar y=\overline{ x+y }\) ainsi que du produit \( k\bar x=\overline{ kx }\) (définition~\ref{PROPooGXMRooTcUGbi}). Toujours par définition du produit appliqué à l'élément \( \bar X\) nous avons \( (\bar X)^2=\overline{ X^2 }\); par récurrence \( \overline{ X^k }=\bar X^k\), et
    \begin{equation}
        \bar Q=\sum_ka_k\bar X^k=Q(\bar X).
    \end{equation}

    Le fait que \( \bar Q=\overline{ Q(X) }\) n'est rien d'autre que le fait que dans \( A[X]\) nous avons \( Q=Q(X)\), comme expliqué dans le lemme~\ref{LEMooGKWQooVOyeDX}.
\end{proof}

%---------------------------------------------------------------------------------------------------------------------------
\subsection{Des sous-modules de polynômes}
%---------------------------------------------------------------------------------------------------------------------------

\begin{normaltext}
Nous allons aussi considérer\nomenclature[A]{\( A_n[X]\)}{les polynômes à coefficients dans \( A\) et de degré inférieur à \( n\)}
\begin{equation}
    A_n[X]=\{ P\in A[X]\tq \deg(P)\leq n \}.
\end{equation}
Cela est un sous-module libre.
\end{normaltext}

%---------------------------------------------------------------------------------------------------------------------------
\subsection{Évaluation}
%---------------------------------------------------------------------------------------------------------------------------

Soit \( P\in A[X]\). À priori, \( P\) n'est qu'une suite dans \( A\) indexée par \( \eN\). 


Nous avons déjà défini son évaluation sur un élément \( \alpha\in A\) dans la définition \ref{DEFooNXKUooLrGeuh} :
\begin{equation}
    P(\alpha)=\sum_ka_k\alpha^k.
\end{equation}
Cette somme est toujours finie.

\begin{normaltext}      \label{NORMooQFTJooLBcPxl}
    L'ensemble \( A[X]\) est une algèbre et donc un espace vectoriel. Il possède un unique élément nul qui est celui dont tous les coefficients sont nuls; cela est immédiat par la construction en tant que suites presque nulles.
\end{normaltext}

Il n'y a a priori pas équivalence entre le fait d'être un polynôme nul et le fait de s'évaluer à zéro sur tous les éléments de \( A\). Cela sera discuté dans le théorème~\ref{ThoLXTooNaUAKR} et l'exemple~\ref{exVQBooBMPLkD}.

\begin{definition}      \label{DEFooRFBFooKCXQsv}
    Soient un anneau \( A\) et un anneau \( B\) qui contient \( A\) (comme sous-anneau). Pour \( \alpha\in B\) nous définissons \( A[\alpha]_B\) comme étant l'intersection de tous les sous-anneaux de \( B\) contenant \( A\).
\end{definition}
Comme dit plus haut, nous nous permettons d'écrire \( A[\alpha]\) sans préciser \( B\) lorsque ce dernier sera clair dans le contexte.

\begin{proposition}     \label{PROPooPMNSooOkHOxJ}
    Soient un anneau \( A\) et un anneau \( B\) qui contient \( A\) (comme sous-anneau). Pour tout \( \alpha\in B\) nous avons
    \begin{equation}
        A[\alpha]=\{ P(\alpha)\tq P\in A[X] \}
    \end{equation}
    où encore une fois, \( P(\alpha)\) est calculé dans \( B\); le contexte est clair là-dessus.
\end{proposition}

\begin{proof}
    Si \( A'\) est un sous-anneau de \( B\) contenant \( A\) et \( \alpha\), alors \( A'\) contient tous les \( P(\alpha)\) avec \( P\in A[X]\). Nous avons donc
    \begin{equation}
        \{ P(\alpha)\tq P\in A[X] \}\subset A[\alpha].
    \end{equation}
    Par ailleurs, \( \{ P(\alpha)\tq P\in A[X] \}\) est un sous-anneau de \( B\) contenant \( A\) et \( \alpha\). Donc \( A[\alpha]\) y est inclus.
\end{proof}

\begin{lemma}[Thème~\ref{THEMEooZYKFooQXhiPD}]       \label{LEMooIDSKooQfkeKp}
    Si \( \eK\) est un corps\footnote{Définition~\ref{DefTMNooKXHUd}.}, alors l'anneau \( \eK[X]\) est euclidien et principal.
\end{lemma}

\begin{proof}
    Vu que \( \eK\) est un corps, tous les éléments sont inversibles et le degré donne un stathme par le théorème~\ref{ThodivEuclPsFexf}. L'anneau \( \eK[X]\) est donc euclidien et par conséquent principal (proposition~\ref{Propkllxnv}).
\end{proof}

Dans le théorème~\ref{ThoCCHkoU} nous donnerons une preuve directe du fait que \( \eK[X]\) est principal en montrant que tous ses idéaux sont principaux. Nous y démontrerons donc un peu moins pour un peu plus cher, mais avec le plaisir de ne pas devoir passer par un stathme.

\begin{definition}[\cite{ooSXFEooEehobn}]  \label{DefDSFooZVbNAX}
    Soit un anneau \( A\). Deux polynômes \( P\) et \( Q\) dans \( A[X]\) sont dits \defe{étrangers}{etranger@étrangers!polynômes} entre eux si \( 1\) est un pgcd\footnote{Définition~\ref{DefrYwbct}.} de \( P\) et \( Q\). Un ensemble de polynômes \( (P_i)_{i\in I}\) est étranger \defe{dans leur ensemble}{étranger!dans leur ensemble} si \( 1\) est un \( \pgcd\) des \( P_i\).

Les polynômes \( P\) et \( Q\) sont \defe{premiers entre eux}{premier!deux polynômes entre eux} si les seuls diviseurs communs de \( P\) et \( Q\) sont les inversibles.
\end{definition}

Les notions de polynômes étrangers entre eux ou de polynômes premiers entre eux ne sont pas identiques, comme le montre l'exemple suivant.

\begin{example}[\cite{MonCerveau}]
    Soient dans \( \eZ[X]\) les polynômes \( P(X)=2X+2\) et \( Q(X)=2X^2+2\). Le nombre \( 2\) est diviseur commun et n'est pas un diviseur de \( 1\). Donc \( 1\) n'est pas un pgcd de \( P\) et \( Q\). Ils ne sont pas étrangers.

    Mais ils sont premiers entre eux parce qu'ils n'ont pas d'autres diviseurs communs que les inversibles (\( 1\) et \( -1\)).
\end{example}

%---------------------------------------------------------------------------------------------------------------------------
\subsection{Polynôme primitif}
%---------------------------------------------------------------------------------------------------------------------------

\begin{definition}\label{DefContenuPolynome}
    Le \defe{contenu}{contenu}\index{polynôme!contenu} du polynôme \( P=\sum_ia_iX^i\in\eK[X]\) est le pgcd de ses coefficients : $c(P)=\pgcd(a_i)$.
\end{definition}

\begin{definition}[Ordre d'un polynôme]
    Soit \( P\) un polynôme irréductible de degré \( n\) sur \( \eF_p[X]\). L'\defe{ordre}{ordre!d'un polynôme} de \( P\) est
    \begin{equation}
        \min\{ k\tq P\divides X^k-1 \}.
    \end{equation}
\end{definition}

\begin{definition}[Polynôme primitif]           \label{DEFooDVOOooKaPZQC}
    Soit \( p\), un nombre premier et \( P\) un polynôme de degré $n$ dans \( \eF_p[X]\). Nous disons que \( P\) est \defe{primitif}{primitif!polynôme} si
    \begin{enumerate}
        \item
            \( P\) est unitaire et irréductible,
        \item
            les racines de \( P\) sont d'ordre \( p^n-1\) dans \( \eF_p[X]/P\).
    \end{enumerate}
\end{definition}

\begin{definition}[Polynôme primitif au sens du pgcd]       \label{DEFooAIYGooRAEfHU}
    Soit un anneau \( A\). Un polynôme \( P\in A[X]\) est \defe{primitif au sens du pgcd}{primitif!polynôme!au sens du pgcd} si ses coefficients sont premiers entre eux.
\end{definition}

\begin{normaltext}
    Pour rappel, il y a plusieurs façons de périphraser le fait que les coefficients soient premiers entre eux. Nous pouvons dire \ldots
    \begin{enumerate}
        \item
            Le pgcd de ses coefficients est \( 1\) parce que c'est la définition~\ref{DEFooXSPFooPumQSy} d'avoir des nombres premiers entre eux.
        \item
            Le contenu de ses coefficients est \( 1\). Parce que le contenu est précisément le pgcd, définition~\ref{DefContenuPolynome}.
    \end{enumerate}
\end{normaltext}

La notion de polynôme primitif au sens du pgcd est particulière aux polynômes à coefficients dans un anneau comme le montre le lemme suivant.

\begin{lemma}
    Si \( \eK\) est un corps, tout polynôme unitaire dans \( \eK[X]\) non nul est primitif au sens du pgcd.
\end{lemma}

\begin{proof}
    Un polynôme unitaire a un \( 1\) parmi ses coefficients, donc le pgcd est forcément \( 1\).
\end{proof}

Lorsque nous utiliserons la notion de polynôme primitif au sens du \( \pgcd\), nous le mentionnerons explicitement. C'est pas exemple le cas pour le corolaire~\ref{CORooZCSOooHQVAOV}.

\begin{definition}[Racine simple et multiple d'un polynôme]
  Soit \( A\) un anneau ainsi qu'un polynôme \( P\in A[X]\) et \( \alpha\in A\) racine de $P$. La \defe{multiplicité}{multiplicité!racine d'un polynôme} de \( \alpha\) par rapport à \( P\) est l'entier \( h\) tel que \( P\) est divisible par \( (X-\alpha)^h\) mais pas divisible par \( (X-\alpha)^{h+1}\).  Nous noterons \( \theta_{\alpha}(P)\)\nomenclature[A]{\( \theta_{\alpha}(P)\)}{la multiplicité de \( \alpha\) par rapport à \( P\)} la multiplicité de \( \alpha\) par rapport à \( P\).
\end{definition}

\begin{normaltext}
    Pour une définition générale d'une racine simple de l'équation \( f(x)=0\), voir la définition~\ref{DEFooXSOQooAnWqKM}. La proposition~\ref{PropHSQooASRbeA} nous indique que toute racine est de multiplicité au moins \( 1\).
\end{normaltext}

\begin{proposition} \label{PropahQQpA}
  L'élément \( \alpha\in A\) est une racine de multiplicité \( h\) du
  polynôme \( P\) si et seulement s'il existe \( Q\in A[X]\) tel que
  \( P=(X-\alpha)^hQ\) avec \( Q(\alpha)\neq 0\).
\end{proposition}

\begin{lemma}[Multiplicités et opérations]       \label{LemIeLhpc}
    Soient \( P\) et \( Q\) des polynômes non nuls de \( A[X]\) et \( \alpha\in A\). Alors
    \begin{enumerate}
        \item
            \( \theta_{\alpha}(P+Q)\leq\min\{
            \theta_{\alpha}(P),\theta_{\alpha}(Q) \}\), et l'égalité a
            lieu si \( \theta_{\alpha}(P)\neq \theta_{\alpha}(Q)\);
        \item     \label{ItemIeLhpciv}
            \( \theta_{\alpha}(PQ)\geq
            \theta_{\alpha}(P)+\theta_{\alpha}(Q)\), et l'égalité a
            lieu si \( A \) est intègre.
    \end{enumerate}
\end{lemma}

Dans le théorème suivant, la partie importante en pratique est la seconde partie parce qu'elle dit que, lorsque nous cherchons les racines d'un polynôme, nous pouvons nous arrêter lorsque nous en avons trouvé autant que le degré, multiplicité comprise.
\begin{theorem} \label{ThoSVZooMpNANi}
    Soit \( A\) un anneau intègre
  et \( P\in A[X]\setminus\{ 0 \}\), un polynôme de degré \( n\). 

  \begin{enumerate}
      \item
  Si \( \alpha_1,\ldots, \alpha_p\in A\) sont des racines deux à deux
  distinctes de multiplicités \( k_1,\ldots, k_p\), alors il existe \(
  Q\in A[X]\), de degré \( n-p\), tel que \(
  P=Q\prod_{i=1}^p(X-\alpha_i)^{k_i}\) et \( Q(\alpha_i)\neq 0\) pour
  tout $i$.
  \item     \label{ITEMooWGOBooGApPOo}
    La somme des multiplicités des racines de \( P\) est au plus \( \deg(P)\).
  \end{enumerate}
\end{theorem}
\index{factorisation!de polynôme}

\begin{proof}
    Si \( p=1\), soit \( \alpha\) une racine de multiplicité \( k\) de \( P\). La définition de la multiplicité d'une racine nous dit que \( P\) est divisible par \( (X-\alpha)^k\) mais pas par \( (X-\alpha)^{k+1}\). Donc il existe \( Q\in \eA[X]\) tel que \( P=Q(X-\alpha)^k\). Il reste à voir que \( Q(\alpha)\neq 0\). Cela est une conséquence de la proposition~\ref{PropHSQooASRbeA} : si \( Q(\alpha)\) était nul, on pourrait lui factoriser \( (X-\alpha)\) et donc avoir \( (X-\alpha)^{k+1}\) qui se factorise dans \( P\), ce qui n'est pas possible.

    Nous supposons que \( p\geq 2\) et nous effectuons une récurrence sur \( p\). Nous considérons donc les \( p-1\) premières racines \( \alpha_1,\ldots, \alpha_{p-1}\) et un polynôme \( R\in\eA[X]\) tel que \( R(\alpha_i)\neq 0\) pour \( i=1,\ldots, p-1\) et
    \begin{equation}
        P=\underbrace{(X-\alpha_1)^{k_1}\ldots (X-\alpha_{p-1})^{k_{p-1}}}_SR.
    \end{equation}
    Par hypothèse \( P(\alpha_p)=S(\alpha_p)R(\alpha_p)=0\). L'anneau \( \eA\) étant intègre, \( S(\alpha_p)\neq 0\) parce que \( \alpha_i\neq \alpha_p\) pour \( i\neq p\). Par conséquent, \( R(\alpha_p)=0\).

    Nous devons encore vérifier que la multiplicité \( \alpha_p\) est \( k_p\) par rapport à \( R\). Pour cela nous utilisons le point~\ref{ItemIeLhpciv} du lemme~\ref{LemIeLhpc} afin de dire que le degré de \( \alpha_p\) pour \( P=SR\) est \( k_p\). Par conséquent
    \begin{equation}
        R=(X-\alpha_p)^{k_p}T
    \end{equation}
    avec \( T(\alpha_p)\neq 0\) et enfin
    \begin{equation}
        P=\prod_{i=1}^p(X-\alpha_i)T.
    \end{equation}
    De plus \( T(\alpha_i)\neq 0\), sinon \( R(\alpha_i)\) serait nul.
\end{proof}

\begin{corollary}       \label{CORooUGJGooBofWLr}
    Un polynôme de degré \( n\) possède au maximum \( n\) racines distinctes.
\end{corollary}

\begin{proof}
    Le théorème \ref{ThoSVZooMpNANi}\ref{ITEMooWGOBooGApPOo} dit que la somme des multiplicités des racines de \( P\) est au maximum \( n\). Mais la proposition \ref{PropHSQooASRbeA} dit que toutes les racines ont une multiplicité au moins un. Donc il ne peut pas y en avoir plus de \( n\).
\end{proof}

\begin{corollary}[Conséquence du lemme de Gauss\cite{ooCDLEooEQGSvn}]       \label{CORooZCSOooHQVAOV}
    Soient \( A\) un anneau factoriel et \( \Frac(A)\) son corps des fractions. Un polynôme non constant \( P\in A[X]\) est irréductible (sur \( A\)) si et seulement s'il est irréductible et primitif au sens du pgcd\footnote{Définition~\ref{DEFooAIYGooRAEfHU}.} sur \( \Frac(A)[X]\).
\end{corollary}

\begin{example}
    Il ne faudrait pas croire qu'être irréductible dans un anneau \( A\) implique d'être irréductible dans le corps des fractions. En effet soit \( A=\eZ[\sqrt{ 5 }]\) et \( P=X^2-X-1\). Nous savons que sa factorisation est 
    \begin{equation}
        P=\left( X-\frac{ 1+\sqrt{ 5 } }{ 2 } \right)\left( X-\frac{ 1-\sqrt{ 5 } }{ 2 } \right).
    \end{equation}
    Si vous ne le saviez pas, faites juste le calcul pour vous en assurer.

    Ce polynôme est irréductible sur \( \eZ[\sqrt{ 5 }]\) mais pas irréductible sur \( \Frac\big( \eZ[\sqrt{ 5 }] \big)\).
\end{example}

%---------------------------------------------------------------------------------------------------------------------------
\subsection{Quelques identités}
%---------------------------------------------------------------------------------------------------------------------------

\begin{lemma}   \label{LemISPooHIKJBU}
    Quelques identités de polynômes.
    \begin{enumerate}
        \item   \label{ItemLTBooAcyMtN}
            Si \( n\) est impair, alors \( 1+X\) divise \( 1+X^n\).
        \item\label{ItemLTBooAcyMtNii}
            Pour tout \( n\) nous avons \( X^n-1=(X-1)(1+X+\cdots +X^{n-1})\).
        \item
            \( X^n-a^n=(X-a)\sum_{i=0}^{n-1}a^iX^{n-1-i}\).
    \end{enumerate}
\end{lemma}

\begin{proof}
  Nous démontrons uniquement le point~\ref{ItemLTBooAcyMtNii}, puisque
  les autres ont été vus en début de chapitre\footnote{Voir l'égalité
    \eqref{Eqarpurmkbk}.}. Le cas \( n=1\) est évident. Procédons
  alors par récurrence en considérant un nombre entier impair \( n\) :
    \begin{subequations}
        \begin{align}
            1+X^{n+2}&=1+X^n+X^{n+2}-X^n\\
                    &=(1+X)P+X^n(X^2-1)\\
                    &=(1+X)P+X^n(X+1)(X-1)\\
                    &=(1+X)\big( P+X^n(X-1) \big).
        \end{align}
    \end{subequations}
\end{proof}
